\chapter{Health Promotion and the Yogic View of the Person}
\label{ch:health-promotion}

\emph{What does it mean to promote health, how does yoga see health and disease, and which attitudes and states of mind does it cultivate?}
This chapter covers the ideas of the health unit that are not about a single system or practice: what health and health promotion are, the holistic view of yoga, the four positive attitudes of the Yoga Sutra, the four bhavas, and the shaping of personality. The common thread is that yoga works on the whole person, mind first, and that its role is prevention and well-being, never cure.

\begin{unverified}
This chapter is written from notes blocks whose fact check is still pending. Points follow the official guidebook (its third, fifth and sixth chapters), the WHO documents and the Yoga Sutra, as marked in the notes; the bhava practices come from the workshop slides only. A reviewer must check it before print.
\end{unverified}

\section{Health and health promotion}

\subsection*{What health is}
The World Health Organization defines health as a state of complete physical, mental and social well-being, and not merely the absence of disease. The guidebook and Indian texts add a fourth, spiritual, side. State both, each with its source. The Sanskrit word \emph{svāsthya} means abiding in one's own self (\emph{sva}, self; \emph{stha}, abiding).

Yoga describes physical well-being as relaxed muscles, loose joints, a calm use of energy, and organ systems that work together even under stress. The texts give early signs of progress in yoga: lightness of body, health, a clear voice and a calm mind.

\subsection*{What health promotion is}
Health promotion, in the WHO's Ottawa Charter of 1986, is the process of enabling people to increase control over their health and to improve it. It is more than treatment: it looks at the whole person and aims at healthy living and the prevention of disease, through healthy habits and awareness.

\subsection*{The part yoga plays}
Yoga builds strength, balance and joint mobility; the cleansing actions cleanse inside; pranayama works on the breath and prana; meditation steadies the mind and attention. It works through conduct, thought, diet, asana, pranayama and meditation together. The guidebook says stiffness and pain can lessen with regular practice. Claims about blood pressure, detoxification or heart function have no source and are not made here.

Yoga is learned by doing. Steady practice (\emph{abhyāsa}, as Patanjali advises) and watching its effects give a better understanding of one's own well-being, and a learner may try several paths, devotion, selfless action or surrender, to see what suits.

\subsection*{The limits of yoga}
Know the contraindications before starting. Yoga is mainly preventive. People differ in constitution, so not every practice suits everyone, and comparing oneself with others is avoided. Yoga works alongside medical care, never instead of it, and it needs the support of a good diet and lifestyle. Say ``promotion and prevention'', and keep medical claims out.

\section{The holistic approach}

\subsection*{The whole person}
Yoga sees the person as more than the body. Its maps of the whole are the five sheaths, the three bodies and the five elements, described in \cref{ch:principles-schools}. The five sheaths, in order, are annamaya (the physical body, made of food), pranamaya (energy), manomaya (mind and senses), vijnanamaya (intellect and discernment) and anandamaya (bliss). The guidebook gives what feeds each sheath: food the annamaya, prana the pranamaya, the drawing-in of the senses (pratyahara) the manomaya, meditation the vijnanamaya, and samadhi the anandamaya.

\subsection*{Mind first, then body}
The texts say that imbalance often begins in the mind, through strong likes and dislikes; it then upsets the flow of prana and finally shows in the body. The mental disturbance is called \emph{ādhi} and the physical illness that follows \emph{vyādhi}; the modern word is psychosomatic. A disorder that starts in the mind (\emph{ādhija vyādhi}) belongs first to the manomaya sheath; one that does not (\emph{anādhija vyādhi}) is mainly of the annamaya sheath.

\subsection*{Prevention}
The guidebook calls yoga basically a preventive discipline. Patanjali says: \emph{heyaṃ duḥkham anāgatam}, the suffering that has not yet come is to be avoided. Disease itself is the first of the nine obstacles to practice in the Yoga Sutra. The holistic aims are physical, mental, social and spiritual well-being, with stress management and a balanced life; yoga's way is to work at every level together.

Interest in breath practices for lung health has grown in recent years, but any such research claim must name its study and stay within well-being; the workshop's claims about managing existing illness are left out.

\section{The four positive attitudes}

The Yoga Sutra (first chapter, sutra 33) names four attitudes, each matched to a kind of person:
\begin{itemize}
\item \textbf{Maitri}, friendliness, towards the happy.
\item \textbf{Karuna}, compassion, towards those who suffer.
\item \textbf{Mudita}, gladness, towards the virtuous and those who do well.
\item \textbf{Upeksha}, even-minded acceptance, neither joining in nor condemning, towards those who act wrongly.
\end{itemize}
The result the Sutra names is \emph{citta prasādana}, a clear and settled mind; the guidebook speaks of peace, calm, serenity and clarity. The attitudes matter because resentment, envy, cruelty and pride disturb the mind. Translations differ: upeksha is rendered as indifference, equanimity or acceptance, and karuna is sometimes read as willingness to help. Learn the four names, the situation each fits, and one plain meaning for each. Teach them as pairs of situation and response; a short role play is more useful than a definition.

\section{The four bhavas}

A bhava is a state of mind: a feeling, emotion or attitude. Yoga practices work on the mind partly by bringing up a particular bhava. The syllabus names four positive bhavas:
\begin{itemize}
\item \textbf{Dharma}: discipline and duty, right conduct.
\item \textbf{Jnana}: knowledge and clarity.
\item \textbf{Vairagya}: detachment, letting go.
\item \textbf{Aishvarya}: inner strength and power.
\end{itemize}
The workshop names four negative bhavas: adharma (unrighteousness), raga (attachment), dvesha (aversion) and ajnana (ignorance). In Samkhya the opposites are given instead as adharma, ajnana, avairagya and anaishvarya; give the workshop list and note the other if asked. The positive bhavas steady the mind and keep conduct, understanding, calm detachment and inner strength; the negative ones are seen as causes of restlessness and distress.

\textit{Source for the practices below: the author's workshop teaching; no Board or guidebook text gives them.} For dharma, steady sitting postures with yama and niyama; for jnana, balancing postures with kapalabhati and trataka; for vairagya, forward-bending postures; for aishvarya, self-awareness and introspection.

\section{Yoga and personality}

\emph{Personality} is a modern word. The texts speak of \emph{svabhāva}, a person's own nature, shaped by the three gunas in different mixtures: sattva (purity, calm, knowledge), rajas (action, drive, restlessness) and tamas (inertia, dullness). The guidebook ties the sattvic mind to calm, alert, kind and thoughtful qualities.

Yoga shapes personality by growing sattva. Right conduct (yama and niyama), a sattvic diet and way of life, asana, pranayama, concentration and meditation make the mind calmer and steadier, and the four attitudes shape conduct towards others. Expected gains of regular practice include self-awareness, emotional control, discipline, confidence, a positive attitude, empathy and resilience; present them as what regular practice tends to bring, not as a guarantee.

\begin{recall}
\begin{itemize}
\item WHO: physical, mental and social well-being; the guidebook adds spiritual. Svasthya: abiding in oneself.
\item Health promotion: enabling people to control and improve their health; prevention, not cure.
\item Mind first: adhi (mental) leads to vyadhi (physical).
\item Attitudes: maitri (happy), karuna (suffering), mudita (virtuous), upeksha (wrongdoers).
\item Bhavas: dharma, jnana, vairagya, aishvarya. Svabhava is shaped by the gunas.
\end{itemize}
\end{recall}

\subsection*{Worked examples}
\begin{itemize}
\item \textbf{A colleague at work is promoted and you feel envious.} Mudita: be glad for the one who does well. Envy disturbs the mind; gladness settles it.
\item \textbf{A student says, ``Yoga cured my friend's diabetes, so I'll stop my tablets.''} Do not agree. Yoga supports prevention and well-being and works alongside medical care; the student must not change any treatment without the doctor.
\item \textbf{Model answer: explain the holistic approach in four lines.} The person is more than the body: five sheaths from food to bliss. Imbalance often starts in the mind (adhi) and then shows in the body (vyadhi). Yoga works at every level together: conduct, diet, asana, pranayama and meditation. Its aim is prevention and well-being, not cure.
\end{itemize}

\subsection*{Check yourself}
\begin{enumerate}
\item How does the WHO define health?
\item What is health promotion?
\item Is yoga a replacement for medical care?
\item Name the five sheaths in order.
\item Where do the texts say imbalance often begins?
\item What is adhija vyadhi?
\item Which attitude fits a person who is suffering? One who acts wrongly?
\item What result does the Sutra name for the four attitudes?
\item Name the four positive bhavas.
\item Which word do the texts use for personality, and what shapes it?
\end{enumerate}

\subsection*{Answers}
\begin{enumerate}
\item A state of complete physical, mental and social well-being, not merely the absence of disease.
\item The process of enabling people to increase control over their health and to improve it.
\item No; it supports well-being and works alongside medical care.
\item Annamaya, pranamaya, manomaya, vijnanamaya, anandamaya.
\item In the mind (the manomaya sheath).
\item A disorder that begins at the mental level.
\item Karuna, compassion; upeksha, even-minded acceptance.
\item A clear and settled mind (chitta prasadana).
\item Dharma, jnana, vairagya and aishvarya.
\item Svabhava, a person's own nature, shaped by the three gunas.
\end{enumerate}

\paragraph{Sources.} WHO Constitution (1948) and the Ottawa Charter for Health Promotion (1986); the official guidebook for Level I (2016, older scheme), third, fifth and sixth chapters; the Yoga Sutra (first and second chapters); Samkhya Karika on the bhavas; the Board syllabus for Level 1.
